% C44_POS /root/scan_w.  New deductions, 2026-11-22 Europe/Berlin.
% The model is GPT-6 Astra Pro.  Runtime reasoning-effort setting is not exposed.
% This file is a paper fragment.  The root paper owns theorem environments.

\section{An algebraic replacement for Wirtinger in the Weil-plane detector argument}

The statements in this section distinguish an unconditional intersection
estimate from the positivity and algebraicity inputs.  Every invocation of
032, Corollaries 8.3 or 8.4, is conditional on the unrefereed preprints
specified by the user.  The characteristic-zero Theorem W in
\texttt{HODGE\_DIVISION\_ATTACK\_20261121.md}, R4, is not being reproved.
We give its missing positive-characteristic degree estimate and then record
the resulting, slightly differently normalized, detector theorem.

\subsection*{C44-W01. A characteristic-free intersection bound on an abelian variety}

\begin{proposition}[C44-W01]
Let $A$ be an abelian variety of dimension $2n$ over an algebraically
closed field, and let $h=c_1(L)$ be an integral ample class.  Put
\[
 D(Z)=\deg(h^n\cdot Z),\qquad
 V=\frac{\deg h^{2n}}{(2n)!}=\chi(L)\in\mathbf Z_{>0},\qquad
 b_h=\frac{2^n}{n!\sqrt V}.
\]
For effective integral cycles $Z,Z'$ of dimension $n$, including cycles
with multiplicities,
\begin{equation}\label{c44:algebraic-bezout}
 0\leq\deg(Z\cdot Z')\leq b_h^2D(Z)D(Z').
\end{equation}
No standard conjecture and no lifting hypothesis is used.
\end{proposition}

\begin{proof}
It suffices by bilinearity to treat integral subvarieties.  On
$Z\times Z'$ consider $f(z,z')=z-z'$ and $s(z,z')=z+z'$.  Translation
moving, or equivalently refined pullback of a general point of $A$ by
$f$, gives
\[
 f_*[Z\times Z']=e[A],\qquad e=\deg(Z\cdot Z')\geq0.
\]
If $f$ is not dominant, both sides have coefficient $e=0$.  If it is
dominant, it is generically finite and $e$ is its degree, including any
inseparable multiplicity.  To see explicitly that the intersection
coefficient agrees with this degree, a general translate of $Z'$ meets
$Z$ properly and avoids the singular loci of both subvarieties: the
bad incidence loci have dimension at most $2n-1$.  At the resulting
zero-dimensional intersections the two smooth subvarieties are local
complete intersections in the smooth variety $A$, and properness gives
the usual length intersection multiplicity.  The same description
applies to the generic fibre of $f$.  Translation preserves cycle
classes.  This also proves the assertion without separability or
transversality assumptions.

Let $h_1,h_2$ be the pullbacks of $h$ to $Z\times Z'$, and set
$\alpha=f^*h$, $\beta=s^*h$.  Both are nef.  The theorem of the cube,
in its parallelogram consequence, gives the numerical identity
\[
 \alpha+\beta=2(h_1+h_2).
\]
Here $[-1]^*h=h$ numerically; no symmetric representative of the
original line bundle is required.  A convenient exact source is
Milne, \emph{Abelian Varieties}, Corollary 6.5, applied to the maps
$z,z',-z'$; the numerical consequence is the displayed identity.
All mixed intersections of nef divisors with an effective cycle are
nonnegative, by approximation by ample rational classes.  Therefore
\begin{align*}
 e\deg h^{2n}
 &=\deg(\alpha^{2n}\cap[Z\times Z'])\\
 &\leq\deg((\alpha+\beta)^{2n}\cap[Z\times Z'])\\
 &=2^{2n}\binom{2n}{n}D(Z)D(Z').
\end{align*}
The last equality uses $\dim Z=\dim Z'=n$.  Division by
$\deg h^{2n}=(2n)!V$ proves the bound.  The asserted formula and
integrality for $V$ are the abelian Riemann--Roch theorem, Milne,
Theorem 13.3(b); ampleness gives $V>0$.
\end{proof}

\paragraph{Realizations and historical scope.}
Equation \eqref{c44:algebraic-bezout} is an inequality between rational
intersection numbers.  It is valid for classes represented in either
$\ell$-adic or crystalline cohomology because the cycle map preserves
these products and trace.  It does not put an order on $\mathbf Q_\ell$.
This is precisely the kind of degree bound needed below.  In
Grothendieck's framework its operations are proper pushforward, refined
intersection, and the biextension identity; SGA 7 I, Exp. VIII,
4.12.1 and 2.3.11 provide the biextension convention already used by
JAC-T4.3.  It supplies a new quantitative input to the positive
incidence calculation of cote 44, scans 008--009, and to Theorem W,
step 3.  It makes no assertion that a new axis class lifts.

\subsection*{C44-W02. The signed Weil-plane Gram matrix in any characteristic}

Fix an imaginary quadratic field $K$ and an order
$\mathcal O\subseteq K\cap\operatorname{End}(A)$ such that the Rosati
involution of $h$ induces conjugation on $K$.  Choose
$a\in\mathcal O$ with $a/\bar a$ not a root of unity.  Write
\[
 \nu=a\bar a\in\mathbf Z_{>0},\quad q=\nu^n,\quad
 s=a^{2n}+\bar a^{2n}\in\mathbf Z,\quad H=4q^2-s^2>0.
\]
For an invertible prime $\ell$, define
\[
 W_\ell=\bigwedge_{K\otimes\mathbf Q_\ell}^{2n}H^1(A,\mathbf Q_\ell)
 \subset H^{2n}(A,\mathbf Q_\ell).
\]
It has dimension two over $\mathbf Q_\ell$.  Let
\[
 f(t)=(t-a^{2n})(t-\bar a^{2n}),\quad
 g(t)=\prod_{j=1}^{2n-1}(t-a^j\bar a^{2n-j}).
\]
Both polynomials have integer coefficients.  Choose $R>0$ and
$u,v\in\mathbf Z[t]$ with $uf+vg=R$, for example the absolute
resultant with the sign of the Bezout relation adjusted.  Put
\[
 p(t)=v(t)g(t)=\sum_jc_jt^j,\qquad
 \kappa=\sum_j|c_j|\nu^{jn},\qquad e_W=R^{-1}p([a]^*).
\]
These are exactly the projector data of the supplied Theorem W.
We use them, with the same normalization, on the new realization.

\begin{proposition}[C44-W02; conditional on unrefereed 032, Corollary 8.4]
Let $Y\subset A$ be an integral $n$-dimensional subvariety with
$e_W[Y]\ne0$, and let $D=D(Y)$.  Then
\[
 y_1=p([a]^*)[Y],\qquad y_2=[a]^*y_1
\]
are independent integral algebraic cycle classes in $W_\ell$.
The form
\[
 Q(x,y)=(-1)^n\deg(xy)
\]
is positive definite on their rational span.  With $d=Q(y_1,y_1)$,
its Gram matrix is exactly
\begin{equation}\label{c44:weil-gram}
 G=d\begin{pmatrix}1&s/2\\s/2&q^2\end{pmatrix},\qquad
 \det G=d^2H/4\in\mathbf Z_{>0},\qquad
 0<d\leq b_h^2\kappa^2D^2.
\end{equation}
Moreover
\[
 W_{\rm alg}:=e_W\operatorname{Alg}^n(A)
    =\mathbf Qy_1\oplus\mathbf Qy_2,
 \qquad W_{\rm alg}\otimes\mathbf Q_\ell=W_\ell.
\]
The two classes have integral signed-cycle representatives whose
positive and negative parts have degrees at most $\kappa D$ and
$q\kappa D$, respectively.
\end{proposition}

\begin{proof}
After extending the coefficient field, write $H^1=V_a\oplus V_{\bar a}$.
Both eigenspaces have dimension $2n$, and the Rosati condition says
that $h$ lies in $V_a\otimes V_{\bar a}$.  Each of the two extreme
exterior powers is therefore annihilated by $h$.  Thus $W_\ell$ is
middle primitive.  The eigenspaces of $[a]^*$ in middle cohomology
have eigenvalues $a^j\bar a^{2n-j}$; all are distinct by the choice
of $a$.  Hence the integral polynomial $p$ is $R$ on the extreme
two lines and zero on the other weight spaces.  This verifies the
projector on this realization, rather than importing a Betti lattice.
The exterior-algebra description is valid integrally for etale
cohomology by Milne, Theorem 15.1.

Every $[a^j]$ is a finite flat isogeny, so the pullbacks in
$p([a]^*)[Y]$ are effective integral cycles with their flat
multiplicities.  Rosati gives $[a^j]^*h=\nu^jh$ numerically and
$\deg[a^j]=\nu^{2nj}$.  Projection formula consequently gives
\[
 D([a^j]^*Y)=\nu^{jn}D.
\]
This proves both asserted degree budgets.  The minimal polynomial
of $[a]^*$ on $W_\ell$ is the irreducible quadratic $f$: indeed
$a^{2n}$ is nonreal.  Hence $y_1,y_2$ are rationally independent.

The signed Hodge standard conjecture, used with exactly the sign
$(-1)^n$, is positive definite on every rational subspace of
primitive algebraic classes.  This is the sole input for positivity
here.  It is the even-degree version of the primitive sign in
cote 44, scan 008; JAC-T4.3 establishes the corresponding incidence
polarization by a curve test.  The unsigned form printed in
Milne's report, Section 3.6, is not used.

Let $T=[a]^*|W_\ell$.  Projection formula gives
$Q(Tx,Ty)=q^2Q(x,y)$.  Thus $T^\dagger=q^2T^{-1}$ and
$T+T^\dagger=s$.  It follows that
$Q(y_1,y_2)=sd/2$ and $Q(y_2,y_2)=q^2d$, proving the matrix
formula.  Its entries are integers, since the $y_i$ are integral
cycle classes, and its determinant is a positive integer.
Expanding $y_1^2$, applying C44-W01 to each pair of effective
pullbacks, and taking absolute values gives
\[
 d=|\deg(y_1^2)|
 \leq b_h^2D^2\Bigl(\sum_j|c_j|\nu^{jn}\Bigr)^2.
\]
This proves the quantitative replacement for Wirtinger.

Finally any finite rationally independent collection of classes in
$W_{\rm alg}$ has a positive definite rational Gram matrix by the
same Hodge-standard input.  Its determinant is nonzero, so those
classes are also $\mathbf Q_\ell$-independent.  There can be at most
two.  The two classes already constructed therefore span
$W_{\rm alg}$ and its realization.
\end{proof}

\paragraph{What ``integral'' means in positive characteristic.}
There is no canonical integral Betti lattice over $\overline{\mathbf F}_p$.
Let $\Lambda$ denote the group of integral algebraic middle classes
modulo numerical equivalence, viewed in cohomology using 032's
numerical--homological equality.  It is a lattice: choose finitely
many integral test cycles giving a nondegenerate rational pairing;
the intersection map embeds $\Lambda$ in a finite-rank free abelian
group.  Write
\[
 \Lambda^\sharp=\{x\in\operatorname{Alg}^n(A):\deg(x\lambda)\in\mathbf Z
                       \text{ for all }\lambda\in\Lambda\}.
\]
The denominator condition needed below is $\delta t\in\Lambda^\sharp$.
The stronger condition $\delta t\in\Lambda$ suffices.  If an integral
compatible realization is specified, its rational classes with integral
pairings also satisfy this condition.  Integrality only at one prime
$\ell$ is insufficient to make the Cramer numerators integers.

\subsection*{C44-W03. A quadratic detector bound in characteristic $p$}

\begin{theorem}[C44-W03; conditional on unrefereed 032, Corollary 8.4]
Retain the preceding notation and suppose $W_{\rm alg}\ne0$.
For $0\ne t\in W_{\rm alg}$ let $\delta\geq1$ be an integer with
$\delta t\in\Lambda^\sharp$, and put $T_0=Q(t,t)>0$.
Define $c_t$ by the same integral effective-cycle class equation as
in the supplied Theorem W, and let $\delta_W$ be its detector degree.
Set $B=b_h\kappa$.  Then
\begin{equation}\label{c44:weil-quadratic}
 \delta_W\leq c_t\leq
 \delta\max\{HB^2,\,4q\kappa B\sqrt{HT_0}\}\,\delta_W^2.
\end{equation}
In particular one can replace $B$ everywhere by the uniform constant
$B_0=2^n\kappa/n!$, independent of the abelian variety, the
characteristic, and the polarization type.  The quartic version
\begin{equation}\label{c44:weil-quartic}
 c_t\leq\delta q^2\max\{B^4,\,4\kappa B^3\sqrt{T_0}\}\,\delta_W^4
\end{equation}
also holds.  Thus the characteristic-$p$ division theorem has both
the requested detector formulation and a quadratic sharpening,
with fully algebraic, explicit constants.
\end{theorem}

\begin{proof}
The hypothesis $W_{\rm alg}\ne0$ supplies a detector: express a
nonzero class as a rational combination of integral effective
cycles and apply $e_W$.  Choose a detector $Y$ of minimal degree
$D=\delta_W$, and form $y_1,y_2$ as in C44-W02.  Put
$b_i=Q(\delta t,y_i)\in\mathbf Z$.  Inverting the exact Gram matrix
\eqref{c44:weil-gram} gives the identity
\begin{equation}\label{c44:exact-clearing}
 \delta Hd\,t=(4q^2b_1-2sb_2)y_1+(4b_2-2sb_1)y_2.
\end{equation}
All coefficients and the positive multiplier $N=\delta Hd$ are
integers.  If $t=t_1y_1+t_2y_2$, Cauchy--Schwarz in this
positive definite rational plane, or the diagonal entries of
$G^{-1}$, gives
\[
 |t_1|\leq\frac{2q\sqrt{T_0}}{\sqrt{dH}},\qquad
 |t_2|\leq\frac{2\sqrt{T_0}}{\sqrt{dH}}.
\]
Therefore the two integral coefficients on the right of
\eqref{c44:exact-clearing} have absolute values at most
$2\delta q\sqrt{HT_0d}$ and $2\delta\sqrt{HT_0d}$.
Using the degree budgets of C44-W02, the positive and negative
parts of the resulting integral cycle each have degree at most
\[
 \kappa D\bigl(2\delta q\sqrt{HT_0d}
                     +2\delta q\sqrt{HT_0d}\bigr)
 \leq4\delta q\kappa B\sqrt{HT_0}\,D^2.
\]
Also $N\leq\delta HB^2D^2$.  These two estimates prove the
upper bound in \eqref{c44:weil-quadratic}.

For completeness the ordinary two-by-two adjugate clearing with
multiplier $\delta\det G$ gives
$N\leq\delta q^2B^4D^4$.  Its two coefficients are bounded by
$2\delta\sqrt{T_0}\|y_1\|\|y_2\|^2$ and
$2\delta\sqrt{T_0}\|y_1\|^2\|y_2\|$, where
$\|y_1\|=\sqrt d\leq BD$ and $\|y_2\|=q\sqrt d$.
Their degree budget is at most
$4\delta q^2\kappa B^3\sqrt{T_0}D^4$, proving
\eqref{c44:weil-quartic}.

Conversely any class equation
$Nt=[Z^+]-[Z^-]$, with $N\geq1$ and effective integral cycles,
has a nonzero Weil projection on at least one integral component
of $Z^+$ or $Z^-$.  That component has degree at most the
corresponding cycle degree.  Hence $\delta_W\leq c_t$.
\end{proof}

\paragraph{Exact change to the supplied Theorem W.}
The original complex constant $|R|/n!$ in the norm estimate is
replaced by $2^n\kappa/(n!\sqrt V)$ for the integral numerator
$y_1$.  We do not claim that Wirtinger's sharper constant survives
unchanged.  All other projector, degree, independence, and integer
clearing steps are now valid without an archimedean realization.
The bound is on cohomological or numerical cycle classes, exactly
as in the supplied definition of $c_t$; it is not an equality in
the Chow group modulo rational equivalence.

\subsection*{C44-W04. The precise balanced-signature hypothesis over a finite field}

\begin{proposition}[C44-W04; Tate-algebraicity conclusion conditional on unrefereed 032]
Let $(A,h,K)$ descend to a finite field $\mathbf F_Q$, with the
$K$-action and polarization defined there, and $\dim A=2n$.
If $p$ splits in $K$, assume the two rational $K\otimes\mathbf Q_p$
components of $A[p^\infty]$ both have dimension $n$ (each has
height $2n$).  If $p$ is inert or ramified, no additional determinant
balance is necessary.  Then Frobenius acts with finite order on
$W_\ell(n)$.  Conditional on 032, Corollary 8.3, the whole
$W_\ell$ is algebraic, $W_{\rm alg}$ is its two-dimensional
rational cycle structure, $\delta_W<\infty$, and C44-W03 applies
to every nonzero class in that rational Weil plane.

When $p$ is split and $p$ does not divide the conductor of
$\mathcal O$, the stated $p$-divisible balance is exactly the usual
$(n,n)$ multiplicity condition on $\operatorname{Lie}A$.
It also holds for the special fibre of a balanced polarized
$\mathcal O_K\otimes\mathbf Z_p$-linear abelian scheme in mixed
characteristic.  These hypotheses give the usual meaningful
all-characteristic version of the signature condition.
\end{proposition}

\begin{proof}
Write $\pi$ for the $Q$-power Frobenius endomorphism.  Its
$K$-linear determinant on $H^1$ is an element $\alpha\in K$,
independent of the realization.  One can see its rationality
without choosing a complex lift: the traces of $b\pi^j$ for
$b\in K$ are rational, by the characteristic polynomial theorem
for endomorphisms, and the nondegenerate trace pairing of $K$
recovers the $K$-linear traces.  Newton identities recover the
$K$-linear characteristic polynomial and its determinant.
The same argument in crystalline cohomology gives the same
determinant: in any of these realizations, the exterior-algebra
description of abelian cohomology and the top-degree action give
$\det(m-[b])=\deg([m]-[b])$ for every integer $m$ for which the
map is an isogeny.  Polynomial identity then makes the
characteristic polynomial independent of the realization.
The trace/characteristic-polynomial input is
Milne, \emph{Abelian Varieties}, Proposition 12.9; it is the
degree-defined characteristic polynomial of an endomorphism.

The determinant $\alpha$ is an algebraic integer, since its
eigenvalues are products of Frobenius eigenvalues.  The
polarization pairing gives
\[
 \alpha\bar\alpha=Q^{2n},\qquad |\alpha|=Q^n
\]
at the two complex embeddings.  The absolute-value assertion is
Weil's theorem for abelian varieties, Milne, Theorem 19.1(c).
In particular $\alpha$ is a unit away from primes over $p$.
If $p$ is nonsplit, its unique prime is fixed by conjugation, so
the norm identity implies
$v(\alpha)=n\,v(Q)$ there.  If $p$ splits, the determinant of
Frobenius on each rational Dieudonne summand has normalized
determinant valuation equal to that summand's dimension:
the Newton polygon of a $p$-divisible group of height $h_0$
and dimension $d_0$ ends at $(h_0,d_0)$, so its slopes sum
to $d_0$.  For exact source statements see Oort,
\emph{Abelian Varieties over Finite Fields}, Sections
21.16--21.19 and Corollary 9.3.  Equivalently each
Dieudonne--Manin factor of dimension $d_0$ and height
$d_0+c_0$ has $d_0+c_0$ slopes $d_0/(d_0+c_0)$.
The balance hypothesis therefore gives
$v(\alpha)=n\,v(Q)$ at both primes.

Consequently $\alpha/Q^n$ is a unit in the imaginary quadratic
field $K$ of absolute value one, hence a root of unity.  On the
two Weil lines the Frobenius eigenvalues are $\alpha$ and
$\bar\alpha$.  A finite extension of the finite field kills
their normalized root-of-unity factors.  The Tate conjecture
in the stated 032 input then gives algebraic classes spanning
$W_\ell$: apply the algebraic projector $e_W$ to cycles spanning
the Tate classes.  C44-W02 identifies their rational cycle
span as a plane.  This proves the assertions.

For split $p$ and a $p$-maximal order, the two integral
idempotents act on the $p$-divisible group and on its Lie
algebra, so the dimensions agree with the two Lie
multiplicities.  In an integral balanced family these ranks
are constant.  Isogenies preserve the dimensions and heights
of the rational $p$-divisible summands.
\end{proof}

\begin{proposition}[C44-W05: a bad-order ambiguity in the word signature]
A determinant identity on $\operatorname{Lie}A$ for an order whose
conductor is divisible by a split prime $p$ does not by itself
imply the balance hypothesis of C44-W04.  In particular the
literal naive such identity cannot replace that hypothesis.
\end{proposition}

\begin{proof}
For a completely concrete example take
$E:y^2=x^3-x$ over $\mathbf F_5$, $K=\mathbf Q(i)$.
It has eight rational points, hence Frobenius trace $-2$;
in particular it is ordinary.  Its order-four automorphism
gives the $\mathbf Z[i]$-action, and its Frobenius is one
of $-1\pm2i$.  Thus $p=5$ splits in $K$, as required.
The following construction also works with any ordinary
CM elliptic curve at a split prime.  After a finite
extension all data are defined over $\mathbf F_Q$.
Let $A=E^{2n}$ with the scalar $K$-action and product polarization,
and restrict the action to $\mathcal O=\mathbf Z+p\mathcal O_K$
(or to $\mathbf Z+p\operatorname{End}(E)$ if necessary).
For $a=m+pb$ the action on the Lie algebra is multiplication by
$m\bmod p$, and the reductions of $a$ and $\bar a$ agree.
Thus
\[
 \det(T-a\mid\operatorname{Lie}A)
 =(T-a)^n(T-\bar a)^n
\]
holds after reduction for every $a\in\mathcal O$, although the
two rational $p$-divisible components have dimensions $2n$ and
$0$.  If $\pi$ is the ordinary Frobenius of $E$, the normalized
Weil determinant is $\pi^{2n}/Q^n$; its valuations at the two
primes over $p$ are nonzero and opposite.  It is not a root of
unity, and the two Weil lines are not Tate.  The example
therefore proves the claimed insufficiency.  It is not a
counterexample to C44-W03 or to the usual balanced PEL meaning
of Weil type.
\end{proof}

\section{Geometric positive transfers and special complete intersections}

The abstract abelian-type transfer remains distinct from a transfer that
preserves the actual incidence form. The following statements give
such geometric transfers and identify the extra special-fibre directions
that they do not reach. Their positivity conclusions are conditional
on the unrefereed 032 input whenever that qualification is displayed.

\subsection*{C44-W06. A finite abelian cover gives an explicit inherited polarization}

\begin{theorem}[C44-W06; conditional on unrefereed 032, Corollary 8.4]
Let $X$ be smooth projective of odd dimension $d=2m-1$ over an
algebraically closed field. Suppose that there is a finite surjective
morphism $f:A\to X$ from an abelian variety. For every rational ample
class $h$ on $X$, put $h_A=f^*h$ and $e=\deg f$. The primitive
incidence quotient $B_X=P^m(X)_{\rm prim}$ has the signed rational
polarization of cote 44, Th.~4. More precisely the induced map
\[
 i_f:B_X\longrightarrow B_A=P^m(A)_{\rm prim}
\]
has finite kernel, and its actual incidence classes satisfy
\begin{equation}\label{c44:finite-cover-polarization}
 i_f^*\ell_{B_A}=e\ell_{B_X},\qquad
 \ell_{B_X}=\frac{(-1)^m}{2}\Delta^*d_{B_X}.
\end{equation}
In particular this is a positive transfer on a concrete class of
nonabelian varieties of geometric abelian type. It does not use
Milne's additional transfer assertion from JAC-P20.
\end{theorem}

\begin{proof}
A finite morphism between smooth varieties of the same dimension is
flat; its degree includes inseparable multiplicity. Push--pull gives
$f_*f^*=e$, and $h_A$ is ample. Lefschetz inverses for $X$ are
algebraic: in degree $i\leq d$, an inverse to $L_h^{d-i}$ is
\[
 \frac1e f_* (L_{h_A}^{d-i})^{-1}f^*.
\]
Indeed $f^*$ commutes with the indicated Lefschetz operators and
the formula composes with $L_h^{d-i}$ to the identity. The inverse
on the abelian variety is algebraic. The same pull--push construction
with its Kunneth projectors supplies those of $X$. Thus the middle
primitive projector is algebraic and self-transpose. Saito's
incidence construction and transpose compatibility, used in
JAC-T4.1, give its nondegenerate factor $B_X$.

Pullback takes primitive cycle families on $X$ to primitive cycle
families on $A$, since $h_Af^*=f^*h$. It therefore induces $i_f$.
Pushforward induces a rational left inverse up to $e$, so $i_f$ has
finite kernel. For two normalized families $z,z'$ on parameter
varieties, projection formula gives the equality of divisorial
incidence correspondences
\[
 (p_{12})_*\bigl((1\times f)^*z\cdot(1\times f)^*z'\bigr)
   =e\,(p_{12})_*(z\cdot z').
\]
By the incidence functoriality of JAC-T4.1, this is the equality
$(i_f\times i_f)^*d_{B_A}=e d_{B_X}$. Diagonal pullback gives
\eqref{c44:finite-cover-polarization}, with the same sign and
normalization as JAC-T4.3. That theorem makes $\ell_{B_A}$ ample
conditional on 032. Pullback by a homomorphism with finite kernel
is ample on its source, and division by $e>0$ preserves rational
ampleness. This proves the conclusion. At the motive level
$e^{-1}f_*f^*=1$ exhibits $h(X)$ as an actual Chow-motive summand
of $h(A)$, so the stated abelian-type scope is also explicit.
\end{proof}

\paragraph{Concrete nonabelian examples.}
In characteristic different from two, take
$A=E_0\times E_1\times\cdots\times E_{d-1}$ with $d\geq3$ odd,
choose a nonzero two-torsion point $a\in E_0$, and quotient by
\[
 (x_0,x_1,x_2,\ldots)\longmapsto(x_0+a,-x_1,x_2,\ldots).
\]
The action is free because of its first coordinate. Its quotient
is smooth projective, the cover is finite etale, and the invariant
$H^1(\mathcal O)$ has dimension $d-1$, so the quotient is not an
abelian variety. In characteristic two use a nonzero three-torsion
translation on $E_0$ and an order-three automorphism of
$E_1:y^2+y=x^3$, acting by $x\mapsto\zeta_3x$. The same argument
works. These examples prove that C44-W06 has nonabelian content in
every characteristic.

\subsection*{C44-W07. The fixed part for every smooth complete intersection}

\begin{proposition}[C44-W07; fixed-part positivity conditional on unrefereed 032]
Let $A$ have dimension $N=2m+c-1$, let $h$ be ample, and let
$k:Y\hookrightarrow A$ be a smooth successive intersection of
$c$ divisors of class $h$, with $\dim Y=n=2m-1$. There is no
generic-last-section hypothesis in this statement. Define
\[
 r=(L_h^c)^{-1}k_*,\qquad
 F=k^*\ker L_h^{c+1},\qquad E=\ker k_*\subset H^n(Y),
\]
and let $e_A$ be the primitive projector on $H^n(A)$. Then
\[
 \Pi_F=k^*e_A r,\quad \Pi_E=1-k^*r,\qquad
 H^n(Y)_{\rm prim}=F\mathbin\perp E
\]
are algebraic self-adjoint projectors and an orthogonal decomposition.
They induce a rational isogeny decomposition of primitive incidence
quotients
\[
 P^m(Y)_{\rm prim}\sim B_F\times B_E.
\]
Its signed incidence class is nondegenerate on both factors and is
a rational polarization on $B_F$, conditional on 032. Consequently
the full Th.~4 assertion for this particular $Y$ is equivalent to
ampleness of its signed incidence class on $B_E$. In particular it
holds whenever $B_E=0$.
\end{proposition}

\begin{proof}
The fixed-and-variable projector proof of H9-T13 uses smooth weak
Lefschetz and the identity $[Y]=h^c$, before any pencil monodromy is
introduced. These hypotheses hold here. Explicitly,
$rk^*=1$, and $k^*r$ is the orthogonal projector onto the ambient
middle cohomology. The primitive condition for an ambient class
$k^*u$ is $L_h^{c+1}u=0$: the Gysin map on $H^{n+2}(Y)$ is
injective by dual weak Lefschetz. The same injectivity makes every
element of $\ker k_*$ primitive. The primitive projector of $A$
is algebraic; the projectors of $Y$ outside its middle degree come
from the weak-Lefschetz identifications with $A$ and their duals.
Thus the displayed middle projectors are actual algebraic
correspondences. Self-adjointness and nondegeneracy give the
incidence decomposition exactly as in JAC-T4.1.

For clarity, the positive scalar in the fixed-part test also remains
unchanged. If a normalized fixed family pulled to a pointed curve
$C$ of genus $b\geq1$ has class $\xi$, put $\eta=(1\otimes r)\xi$
and lift its parameter degree-one component to an algebraic class
$w$ on $J(C)\times A$. Then $h^{c+1}w=0$ and, for
$H=\theta+h$,
\[
 H^{b+c}w=0,\qquad
 \int_{J(C)\times A}w^2H^{b+c-1}
   =\frac{(b+c-1)!}{c!}\int_{C\times Y}\xi^2.
\]
This is H9-T13, equation (H9.47); its proof needs no genericity.
The coefficient is positive. Conditional on 032, the signed curve
degree is nonnegative, with equality only for a zero mixed class.
The nondegenerate signed divisor is therefore ample by the
nef-plus-nondegenerate criterion of JAC-T4.3. Orthogonality makes
ampleness of the full signed divisor equivalent to ampleness of
both factors, proving the final assertions.
\end{proof}

\subsection*{C44-W08. Specialization of the inherited incidence polarization}

\begin{theorem}[C44-W08; conditional on unrefereed 032]
Let $\mathscr Y\to U$ be the smooth locus of a genuine Lefschetz
pencil in the setting of H9-T13, and let $s\in U$ be any geometric
point. Write $B_{\bar\eta}=P^m(Y_{\bar\eta})_{\rm prim}$ for the
geometric generic primitive incidence quotient. After a finite
extension of the local generic field, $B_{\bar\eta}$ has good
reduction at $s$, and its special fibre maps with finite kernel to
$P^m(Y_s)_{\rm prim}$. The actual signed incidence form is positive
on this inherited image. Hence Th.~4 holds for $Y_s$ whenever
\begin{equation}\label{c44:no-incidence-jump}
 \dim P^m(Y_s)_{\rm prim}=\dim P^m(Y_{\bar\eta})_{\rm prim}.
\end{equation}
More particularly, if the generic variable incidence factor of the
pencil is nonzero, then Th.~4 holds for every smooth geometric fibre
of the pencil. If that generic factor is zero, only the new
variable incidence directions that occur at a special fibre remain
outside this theorem.
\end{theorem}

\begin{proof}
Work over the strict henselian local trait at $s$, or its completion,
and pass to a finite extension of its fraction field over which an
abelian representative $B$ and a normalized rational Poincare cycle
$z$ are defined. These are finite data; rational descent is the one
fixed in JAC-P00. As in H9-T5, equations (H9.10)--(H9.11), the
mixed class of $z$ injects the Tate module of $B$ into the appropriate
Tate twist of $H^n(Y_{\bar\eta})$. Indeed its self-incidence is
nondegenerate, so a vector killed by the family map would be in
the radical of that incidence pairing. Smooth proper base change
makes the target unramified at this smooth fibre. Thus $B$ has
good reduction by the Neron--Ogg--Shafarevich criterion, just as in
H9-T5. One exact primary author source for the criterion is Milne,
\emph{Complex Multiplication}, Theorem 6.12, p.~51, which cites
Serre--Tate (1968), Theorem 1.

Let $\mathscr B$ denote the resulting abelian scheme over the trait.
Take the closure of a rational-cycle representative of $z$ in
$\mathscr B\times\mathscr Y$. These schemes are smooth over a
regular trait, so the refined fibre restrictions and intersection
products are defined. Its special-fibre mixed class is the
specialization of the generic mixed class. In particular it is
still in parameter degree one and is primitive on $Y_s$; these
conditions are kernels of maps of the smooth-proper cohomology
systems. Rigidification or degree-one projection, if desired,
alters no incidence class.

Form the actual relative self-incidence divisor by the same proper
pushforward of the product of the two cycle families. Clear its
rational denominator. Its signed half-diagonal pullback is a line
bundle $\mathscr L$ on $\mathscr B$, and refined base change says
that its two fibre classes are exactly the corresponding signed
incidence classes, up to the one positive clearing factor. The
generic class is ample by H9-T13.

Here is the specialization step without an assertion that arbitrary
special cycles lift. Choose a relatively ample line bundle
$\mathscr H$ on $\mathscr B$. Flatness and the Hilbert polynomial
make
\[
 \chi(L_t+uH_t)=\frac{(c_1(L_t)+u c_1(H_t))^{\dim B}}{(\dim B)!}
\]
the same real polynomial in $u$ for the generic and special fibres
$t$. On the generic fibre this polynomial is positive for every
$u\geq0$, since $L_{\eta}+uH_{\eta}$ is ample. Thus on the
special fibre $\chi(L_s)\ne0$ and the polynomial has no positive
real root. Mumford's index theorem gives index zero for $L_s$;
the nondegenerate effective criterion makes it ample. These are
exactly the two abelian-variety criteria used in the checked proof
of JAC-T4.3 (Mumford, \emph{Abelian Varieties}, Chapter III,
Section 16, and Chapter II, Section 6, Application 1).

The specialized family induces a homomorphism
$u:B_s\to P^m(Y_s)_{\rm prim}$. Saito's transpose and incidence
compatibility identifies $u^*\ell_s$ with the ample class just
constructed. Therefore $u$ has finite kernel: restriction to any
positive-dimensional connected kernel would make an ample class
zero. This proves positivity on the inherited image. Under
\eqref{c44:no-incidence-jump} the map is an isogeny, and ampleness
descends through that isogeny.

For the final assertion apply the same construction to the generic
variable factor alone. If that factor is nonzero, the irreducibility
after finite-index restriction in H9-T13, or H9-T5 (H9.11), forces
its family map to fill the entire variable local system $E$.
Specialization preserves this rank. The special variable incidence
factor injects into $E_s$ by its own nondegenerate self-incidence,
so its dimension is at most $\operatorname{rank}(E)/2$. The
inherited image already has that dimension. It therefore maps
isogenously onto the special variable incidence factor and gives
its positive signed form. C44-W07 supplies the fixed factor for
every smooth fibre. This proves the stated all-smooth-fibre scope.
\end{proof}

\paragraph{The remaining special-fibre gate.}
When the generic variable incidence quotient vanishes, the
construction specializes only the fixed incidence quotient. A
special fibre may have a new variable incidence factor. Neither
smooth proper base change nor constancy of the rank of $E$ makes
that factor extend as algebraic cycle families through the pencil.
Its signed positivity is the smaller unresolved statement. Thus
C44-W08 does not turn H9-T13 into an unconditional assertion about
every smooth complete intersection; it proves precisely the
rank-preserving specializations and the nonzero-generic-variable
case, with 032 retained as the positivity input.

\subsection*{C44-W09. The full geometric-abelian-type extension with its extra input exposed}

\begin{proposition}[C44-W09; conditional on unrefereed 032 and the separate JAC-P20 transfer]
Assume, in addition to 032, Milne's Proposition 4.2 in
\emph{Grothendieck's standard conjecture of Lefschetz type over
finite fields}, version 14 July 2022, in the total-cohomology form
specified and audited in JAC-P20. Then JAC-T4.3 holds for every
smooth projective odd-dimensional variety of geometric abelian
type, with its actual signed primitive incidence divisor.
\end{proposition}

\begin{proof}
JAC-P16 gives algebraic Lefschetz and primitive projectors. For
every pointed test curve $C$, the product $J(C)\times X$ is still
of geometric abelian type. JAC-P20's actual abelian cohomology
injection and the additional transfer assertion give the Hodge
standard conjecture for this product. The construction of
JAC-T4.2, using the algebraic Abel inverse on the first factor,
now gives an algebraic primitive class $w$ on $J(C)\times X$ and
the exact identity
\[
 \int_{J(C)\times X}w^2(\theta+h)^{b-1}
       =(b-1)!\int_{C\times X}\xi^2.
\]
The Hodge-standard input on this particular product makes its
signed value nonnegative. Saito's nondegenerate incidence factor
and JAC-T4.3's curve criterion then make the signed divisor ample.
This argument proves the stated conditional corollary; it does
not establish the additional transfer assertion itself.
\end{proof}

\paragraph{Framework and precise dependency.}
C44-W06 proves a positive scalar comparison of the actual forms,
not merely an embedding of motives. C44-W07--W08 use smooth
proper base change, SGA 7's fixed and variable projectors, and
the extension of actual cycle correspondences across a regular
trait. Their target is cote 44, scans 008--009, Th.~4 and its
primitive incidence higher Picard variety. The new residual gate
is the positive sign on a special new variable incidence factor,
not the already checked generic pencil sign. For all varieties
of geometric abelian type, JAC-P20 still leaves the comparison
between an abstract transported positive form and the particular
geometric form as an additional dependency. C44-W09 records
the exact implication if that dependency is supplied.


\section{The supersingular Fermat test for cote 16}

\subsection*{C44-S01. The proposed Fermat candidate already has the required sign}

\begin{proposition}[C44-S01; unconditional]
Let $X=X_N^{2r}\subset\mathbf P^{2r+1}$ be the smooth Fermat
hypersurface of degree $N\geq3$ in characteristic $p\nmid N$,
with $r\geq2$.  Suppose $p^\nu\equiv-1\pmod N$ for some
$\nu\geq1$.  Then all even cohomology is algebraic and all odd
cohomology vanishes.  Numerical and homological equivalence
coincide.  The signed form
\[
 Q_V(x,y)=(-1)^r\deg(xy),\qquad
 V=\ker(h:H^{2r}(X)\longrightarrow H^{2r+2}(X)),
\]
is positive definite on its rational algebraic cycle structure.
In particular the Hodge standard conjecture for $X$ is known
unconditionally; it is not a consequence that needs 032.
\end{proposition}

\begin{proof}
Put $q=p^\nu$ and $b=(q+1)/N$.  The Hermitian Fermat
hypersurface $H=X_{q+1}^{2r}$ maps finitely and surjectively to
$X$ by
\[
 [x_0:\cdots:x_{2r+1}]
 \longmapsto[x_0^b:\cdots:x_{2r+1}^b].
\]
Its degree is $b^{2r+1}$: $f^*h_X=bh_H$ and
$\int_Hh_H^{2r}=q+1=bN$ determine the degree by projection
formula.  In particular $f_*f^*=b^{2r+1}$ and $f^*$ is injective.
Shimada, \emph{Lattices of Algebraic Cycles on Fermat Varieties
in Positive Characteristics}, Theorem 1.1(1)--(2) and
Corollary 1.2, proves that linear cycles span the middle
cohomology of $H$ and that its signed primitive intersection
form is positive definite.  Weak Lefschetz gives the other
cohomology groups.  Pushing forward these algebraic classes
shows algebraicity of all cohomology of $X$.

If $x\in V$, then
$h_Hf^*x=b^{-1}f^*(h_Xx)=0$, and
\[
 (-1)^r\int_H(f^*x)^2
  =b^{2r+1}(-1)^r\int_Xx^2.
\]
The injectivity of $f^*$ and Shimada's positive form prove the
claim.  In all other degrees $X$ has the ambient projective
space Lefschetz strings, whose sign is immediate.  The
intersection pairing is nondegenerate on the rational
algebraic classes because these span cohomology; thus
numerical and homological equivalence agree.  Since
$\operatorname{NS}(X)_\mathbf Q=\mathbf Qh$, this also covers
every ample polarization.
\end{proof}

\paragraph{Exact status of the proposed test.}
Shioda--Katsura, \emph{On Fermat Varieties}, Theorem II,
printed p.~97, says for even dimension and $N\geq4$ that
algebraic generation of cohomology is equivalent to the
displayed congruence.  Thus, in that range, the entire stated
supersingular Fermat candidate class is covered by C44-S01.
The new argument above is only the finite-cover descent
of Shimada's already known sign.  It neither attributes a
new theorem to Grothendieck nor treats algebraicity alone as
a positivity theorem.  The square needed in scan 049 is
handled explicitly below.

\subsection*{C44-S02. An explicit smooth Segre hyperplane}

\begin{proposition}[C44-S02; unconditional]
Let $X=X_N^d$ be a smooth Fermat hypersurface with $N\geq3$.
Choose a finite field $k_0$ containing every $(N-2)$-nd root
of unity in $\overline{\mathbf F}_p$.  Let $M>N(d+1)$ and
choose $\beta$ of degree $M$ over $k_0$ in its extension of
degree $M$.  Define $c_i=\beta^{(N-2)i}$ for
$0\leq i\leq d+1$.  Then
\[
 W=\left\{\sum_{i=0}^{d+1}c_ix_iy_i=0\right\}
       \subset X\times X
\]
is a smooth hyperplane section for the Segre polarization
$L=h_1+h_2$.  Also $Y=X\cap\{x_{d+1}=0\}$ is a smooth
hyperplane section of $X$.
\end{proposition}

\begin{proof}
The assertion for $Y$ follows from $p\nmid N$.  If $W$ had
a singular point, the differential of its equation would
annihilate both projective tangent spaces.  Since the
Fermat gradients are $N(x_i^{N-1})$ and $N(y_i^{N-1})$,
there would be scalars $\lambda,\mu$ with
\[
 c_iy_i=\lambda x_i^{N-1},\qquad
 c_ix_i=\mu y_i^{N-1}\quad\text{for every }i.
\]
Both scalars are nonzero, since all $c_i$ are nonzero and
neither projective coordinate vector is zero.  The nonzero
coordinates of $x$ and $y$ have the same nonempty support
$S$.  Substitution gives, for $i\in S$,
\[
 (x_i^N)^{N-2}=(\mu\lambda^{N-1})^{-1}c_i^N.
\]
Choose $\rho$ with $\rho^{N-2}=(\mu\lambda^{N-1})^{-1}$.
It follows that
$x_i^N=\rho\zeta_i\beta^{Ni}$ for roots of unity
$\zeta_i\in k_0$ with $\zeta_i^{N-2}=1$.
This remains valid if $p\mid(N-2)$, since only the distinct
roots are used.  The Fermat equation for $x$ now gives
\[
 \sum_{i\in S}\zeta_i\beta^{Ni}=0.
\]
This is a nonzero polynomial over $k_0$ of degree at most
$N(d+1)<M$ vanishing at $\beta$, a contradiction.  The
Jacobian criterion therefore proves smoothness.
\end{proof}

\subsection*{C44-S03. The exact square sign and all three corrected scan-050 conditions}

\begin{theorem}[C44-S03; unconditional]
For $X=X_N^{2r}$ of C44-S01, put $d=2r$, $T=X\times X$,
$L=h_1+h_2$, and take $Y,W$ as in C44-S02.  Then
\[
 C_{\rm G}(X),\quad C_{\rm G}(Y),\quad R_0(T,W),\quad
 G_{L,0}(T,W),\quad A_0(T),\quad G_0(T,W),\quad A_0(W)
\]
all hold with the meanings fixed by SIGN\_NEXT SN-C03--SN-C04.
The Hodge standard conjecture $I(T)$ holds.  In its middle
degree the primitive space and form are explicitly
\begin{gather}
 P^{2d}(T)=\mathbf Qw\ \perp\ (V\otimes V),\qquad
 w=\sum_{i=0}^{d}(-1)^ih_1^ih_2^{d-i},\label{c44:fermat-square-primitive}\\
 \int_T(cw+z)^2=(d+1)N^2c^2+(Q_V\otimes Q_V)(z,z)>0
 \quad\text{if }(c,z)\ne(0,0).\label{c44:fermat-square-sign}
\end{gather}
The ring of degree-zero rational correspondence classes
modulo homological equivalence is
\[
 \prod_{i=0}^{d}\operatorname{End}_{\mathbf Q}
       \bigl(\operatorname{Alg}^i(X)\bigr)
 \simeq \mathbf Q^d\times M_{b_d(X)}(\mathbf Q),
\]
and is therefore semisimple.  These are the actual sign and
semisimplicity conclusions of cote 16, scan 049, for this
nonabelian family.
\end{theorem}

\begin{proof}
Weak Lefschetz gives the orthogonal decomposition
\[
 H^{\rm even}(X)=U\perp V,\qquad
 U=\mathbf Q\langle1,h,\ldots,h^d\rangle,
\]
where $V$ is in cohomological degree $d$ and $hV=0$.
The rational structures here are given by cycle classes;
they become the full cohomology after scalar extension.
The middle cohomology of $T$ consequently consists of the
ambient terms, the two terms $V\otimes h^r$ and
$h^r\otimes V$, and $V\otimes V$.  On the ambient
terms the kernel of multiplication by $L$ is the single
line displayed in \eqref{c44:fermat-square-primitive}:
successive coefficients satisfy $c_i+c_{i-1}=0$.
Multiplication by $L$ is injective on each of the two
cross terms and their images have different Kunneth
degrees.  It is zero on $V\otimes V$.  Orthogonality
follows from $\int_Xvh^r=0$ for $v\in V$.

In $w^2$ the $d+1$ contributing pairs have indices
$i,d-i$ and sign $(-1)^d=1$, giving
$\int_Tw^2=(d+1)N^2$.  The pairing on $V\otimes V$
is the tensor product of the two raw pairings; their
two factors $(-1)^r$ cancel.  It is therefore
$Q_V\otimes Q_V$, a positive definite form.  This proves
the exact critical sign.

For the full Hodge standard statement, decompose
$H^*(T)$ into the four $L$-stable blocks
$U\otimes U$, $V\otimes U$, $U\otimes V$, and
$V\otimes V$.  They are orthogonal.  The first has
the Lefschetz module and intersection forms of
$\mathbf P^d\times\mathbf P^d$, multiplied by $N^2$;
its signed forms are positive by characteristic-zero
Hodge--Riemann, with identical rational matrices in
every characteristic.  In $V\otimes U$ the primitive
classes are exactly the base $V\otimes1$, in
codimension $r$.  Their signed form is $NQ_V$;
the same is true for $U\otimes V$.  The last block
is primitive in middle degree and was just treated.
This proves $I(T)$ in all degrees.  Every ample
class on $T$ is a positive linear combination of
$h_1,h_2$; the same argument, with positive rescalings,
also handles these polarizations if required.

For a smooth hypersurface $Z$ of degree $N$ and
dimension $e$, the Lefschetz inverses below its middle
degree are already explicit.  For even $i<e$ the
correspondence
\[
 \frac1N h_Z^{i/2}\times h_Z^{i/2}
\]
induces the inverse of
$L_Z^{e-i}:H^i(Z)\longrightarrow H^{2e-i}(Z)$;
for odd $i<e$ both spaces vanish.  The middle inverse
is the identity.  Thus $C_{\rm G}(X)=B(X)$ and
$C_{\rm G}(Y)=B(Y)$.  The product stability of $B$
(Kleiman 1968, Proposition 2.12), or the cycle-basis
matrix units below, gives $B(T)$ and in particular
the algebraicity of $(L^2)^{-1}$ in the critical
degrees of SN-C03.

Write $j:W\hookrightarrow T$.  If $j^*x$ is algebraic
for $x\in H^{2d-2}(T)$, then
$L j_*j^*x=L^2x$ is algebraic; applying the algebraic
inverse proves $R_0(T,W)$.  For
$\gamma\in\operatorname{Alg}^{d+1}(T)$, set
$x=(L^2)^{-1}\gamma$ and $y=j^*x$.  Both are
algebraic and $Lj_*y=\gamma$, proving $G_{L,0}$.
The critical $A_0(T)$ was used here; the equivalence
in SN-C03 supplies $G_0$ and $A_0(W)$ as well.
This verifies the maps with the required extra $L$;
no unweighted lower-Gysin replacement is made.

Finally choose rational cycle bases for each
$\operatorname{Alg}^i(X)$ and the intersection-dual
bases in degree $d-i$.  Exterior products of these
classes give all rational matrix units on
$H^{2i}(X)$.  Every degree-zero correspondence
acts degreewise; Kunneth and Poincare duality make
this action faithful modulo homological equivalence.
It is therefore exactly the product of the stated
full matrix algebras.  Since all nonmiddle Betti
numbers in even degree are one, the displayed
description follows.
\end{proof}

\paragraph{What task 4 does and does not close.}
C44-S01--C44-S03 rule out the suggested supersingular Fermat
family as an example whose square-sign status is unknown.
They verify the corrected scan-050 proposition and the
scan-049 sign and correspondence ring explicitly and
unconditionally on that family.  A different smooth
projective $X$ with a genuinely unestablished $I(X\times X)$
is still required for the adversarial part of task 4; absence
of a source for one named equation would not certify that
status.  None of these findings closes the special-axis
sign for a complete intersection in an abelian variety.

\paragraph{Framework and programme crosswalk.}
For C44-S01--S03 the relevant archive locators are
\texttt{ARCHIVE\_16\_scan049.tex}, the semisimple
correspondence paragraph, and
\texttt{ARCHIVE\_16\_scan050.tex}, the proposition and
final weighted diagram.  They instantiate SIGN\_NEXT
SN-C03--SN-C04 and its successor question 4.  Their
Lefschetz inverses also give the $C_{\rm G}$ input to
cote 44, Th.~6, scan 004, through JAC-P02.  Weak
Lefschetz and Gysin are the Weil-cohomology operations
in SGA 5, Exp. VII, Section 7 (complete intersections).
The finite power maps and explicit hyperplanes are
scheme morphisms over finite extensions of $\mathbf F_p$;
the identities survive every Weil cohomology realization.
At the level of homological motives the Fermat object
is a direct sum of Tate objects with the positive
Lefschetz forms just computed.  No equality with its
Chow motive modulo rational equivalence is silently
deduced.  C44-W01--W05 use the exterior-power, hence
$\lambda$-operation, of $h^1(A)$, the algebraic
endomorphism projector, and rational cycle intersections.
No additional $K$-theory assertion, Selmer injection,
or motivic $L$-function nonvanishing is inferred.
